\section{Axis trials: is the boat moving the right way?}
\label{sec:trial}

An axis trial holds neutral for a baseline window, then a command on one axis,
then neutral again. Let $c$ be the velocity component for the axis
($v_x$, $v_y$ or $\omega_z$), $B$ the baseline samples and $D$ the driven samples.
\[
  \bar c_B = \frac1{|B|}\sum_{k\in B} c_k,\qquad
  \bar c_D = \frac1{|D|}\sum_{k\in D} c_k,\qquad
  r = \bar c_D - \bar c_B .
\]
The \emph{response} $r$ subtracts the drift already present before the command,
so a current pushing the boat forward does not make a reverse command look
weak. The verdict compares its sign with the commanded direction $s=\pm1$:
\[
  \text{verdict} = \begin{cases}
    \text{LOW RESPONSE} & |r| < 0.03\\
    \text{SIGN MATCH} & r\,s > 0\\
    \text{SIGN MISMATCH} & \text{otherwise.}
  \end{cases}
\]
\paragraph{When not to trust it.} The trial is INCONCLUSIVE with fewer than five
driven or three baseline samples, with a gap between samples longer than the
staleness limit, or if the odometry frames change mid-trial. A sample stamped at
the exact instant the command starts is not counted as driven: it cannot reflect
the command yet. SIGN MATCH says the wiring and signs are right; it is not a
calibration (that is Section~\ref{sec:char}).
